\documentclass[10pt,a4paper]{article}
\usepackage[margin=25mm]{geometry}
\usepackage{amsmath,amssymb,amsthm,booktabs,array,longtable}
\usepackage[T1]{fontenc}
\usepackage{lmodern,microtype,xurl}
\usepackage[hidelinks]{hyperref}
\setlength{\parskip}{4pt}
\setlength{\parindent}{1em}
\setlength{\emergencystretch}{2em}
\renewcommand{\arraystretch}{1.13}
\renewcommand{\topfraction}{.9}
\renewcommand{\bottomfraction}{.6}
\renewcommand{\textfraction}{.08}
\renewcommand{\floatpagefraction}{.75}
\setcounter{topnumber}{3}
\setcounter{totalnumber}{4}
\newtheorem{theorem}{Theorem}
\newtheorem{lemma}[theorem]{Lemma}
\newtheorem{proposition}[theorem]{Proposition}
\theoremstyle{remark}
\newtheorem{remark}[theorem]{Remark}
\newcommand{\R}{\mathbb R}
\newcommand{\norm}[1]{\left\lVert #1\right\rVert_\infty}
\newcommand{\abs}[1]{\left|#1\right|}
\newcommand{\Cert}{\mathrm{C}}
\newcommand{\Uncert}{\mathrm{U}}
\newcommand{\Partial}{\mathrm{P}}
\newcommand{\Notrun}{\mathrm{N}}
\newcommand{\Hold}{\mathrm{H}}
\title{Interval certification of selected equal-mass\newline
three-body periodic brake orbits}
\author{Tatsuhiko Takimoto\\\small Independent researcher}
\date{Version 1.0.0, 25 October 2026}
\begin{document}
\maketitle
\begin{abstract}
We give interval-arithmetic existence certifications for parameter boxes
associated with 13 of the first 20 rows of a published table of equal-mass
planar three-body periodic brake orbits. An Arb implementation certifies
rows 1--6 and 10--16. Separate flow computations certify ten of these
rows: a handwritten MPFR integrator certifies rows 1, 2, 3 and 12, and a
CAPD multiprecision C1-Lohner implementation certifies rows 4, 5, 6, 10,
13 and 14. Its six selected paths have source and saved-trace audits
of the local inclusions and their connection to the original family.
Equal centers, nested parameter boxes and uniqueness identify the same
normalized solution in each pair. The argument combines quarter-period
symmetry, validated flow and derivative enclosures, a three-variable
Krawczyk contraction and positive pairwise-distance bounds. The CAPD
route uses nonuniform time steps and forward elimination of the
multiple-shooting connection derivatives. Each certificate gives a
unique reduced zero in its stated box and a nonconstant, collision-free
periodic solution. The implementations share the model and symmetry
proof. Row counts do not establish geometric distinctness; least period,
stability and catalogue completeness are outside the claim. Seven entries
numerically match previously catalogued self-dual candidates; earlier
labels for rows 4, 10 and 12 were not identified in the tables reviewed.
We make no claim of priority.
\end{abstract}

\section{The question and the result}
High-precision numerical data locate a candidate periodic orbit, but a
small return residual alone does not establish the existence of an exact
periodic solution. We study this question for the Euler-coordinate data
of Hristov et al.\ \cite{Hristov,Data}. The concrete object to be certified
is a parameter box around the finite decimals in a table row. The
numerical discovery of these candidates belongs to the catalogue authors.
A zero of
a reduced quarter-period equation supplies the exact parameters; a
symmetry argument then supplies the full periodic solution.

\begin{theorem}[Computer-assisted result]
\label{thm:result}
Let $x_{0,j}=(a_{0,j},b_{0,j},T_{0,j})$ be the exact rational triple given
by the first three decimal fields of row $j$ of
\texttt{euler\_100\_4860.txt}, with row numbers starting at one. For
\[
 \mathcal J=\{1,2,3,4,5,6,10,11,12,13,14,15,16\},
\]
the Arb computations certify exactly one zero $x_j=(a_j,b_j,T_j)$ of
the map $F$ defined in Section~\ref{sec:model} in
\[
 X_j=x_{0,j}+[-r_j,r_j]^3,\qquad
 r_j=\begin{cases}2^{-120},&j=2,3,\\2^{-160},&j\in\mathcal J\setminus\{2,3\}.
 \end{cases}
\]
The solution with these parameters is nonconstant, is collision-free for
all real time, and has $T_j>0$ as a period. Its pairwise distances have
the positive lower bounds in Table~\ref{tab:distances}.

For $j\in\{1,2,3,4,5,6,10,12,13,14\}$, a second flow implementation also
certifies exactly one zero in $Y_j=x_{0,j}+[-s_j,s_j]^3$, where
\[
 s_j=\begin{cases}
 2^{-90},&j=4,13,\\
 2^{-100},&j=1,2,3,12,14,\\
 2^{-130},&j=6,\\
 2^{-140},&j=5,10.
 \end{cases}
\]
The second computation uses the handwritten MPFR integrator for
$j=1,2,3,12$ and CAPD for $j=4,5,6,10,13,14$. In each case its zero is $x_j$.
\end{theorem}

The uniqueness assertion concerns $F$ with labels, Euler initial
configuration, orientation and scale fixed. It is not uniqueness in the
space of all periodic three-body motions. A certified $T_j$ need not be
the least period. Thus the theorem concerns 13 row-associated boxes,
without counting geometrically distinct orbits. The numerical part of
the proof is the set of validated flow and derivative enclosures. The
analytic implications of those enclosures are proved below.

\subsection{Scope of assurance}
\label{sec:assurance}
Table~\ref{tab:assurance} gives the current guarantee structure. The
three implementations use a common model, reduced map and prose
symmetry proof. Arb and handwritten MPFR share the Taylor sufficient
conditions below; CAPD supplies a different flow-enclosure route.
Code separation therefore tests implementation agreement under common
mathematical premises. The MPFR and CAPD routes both use MPFR scalar
arithmetic. This is a computer-assisted prose argument with the trusted
numerical components stated in the table.

\begin{table}[htbp]
\centering\small
\begin{tabular}{@{}p{.25\linewidth}p{.69\linewidth}@{}}
\toprule
Component & Present assurance and extent\\
\midrule
Quarter-period reduction & Prose proof in Lemma~\ref{lem:symmetry}; the
lemma is a common premise of all three implementations.\\
Flow and derivative & Arb validates the 13 stated rows; MPFR separately
validates four; CAPD validates six further rows. Outward arithmetic,
the respective integrator, compiler and runtime belong to the trusted
computing base.\\
Terminal algebra & Exact rational checks of the saved intervals verify
the finite matrix and contraction tests. They take the saved flow and
derivative enclosures as inputs.\\
Agreement & Equal rational centers, nested boxes and uniqueness identify
the same reduced zero for rows 1, 2, 3, 4, 5, 6, 10, 12, 13 and 14.\\
CAPD audit & Separate integer/rational checks of the saved local
inclusions, time choices, derivative reduction and terminal data cover
the six selected paths. Taylor coefficients and set updates retain a
fixed-source trust premise; these audits are not new ODE integrations.\\
Formal and human review & No Lean formalization and no human expert
review of this work. The scope and outcome of AI mathematical reading
are recorded in the accompanying verification record.\\
Claim boundary & Existence, a positive collision margin and local
parameter-box uniqueness; least period, stability, geometric
distinctness and completeness of the table remain open here.\\
\bottomrule
\end{tabular}
\caption{The layers of the computer-assisted argument.}
\label{tab:assurance}
\end{table}

The audits do not independently regenerate Taylor coefficients, rebuild
the cloud binaries, or formally verify the software. Compiler,
MPFR/GMP, CAPD arithmetic, automatic differentiation and set operations,
hardware, and the correspondence between fixed sources and recorded
executions remain trusted. The methods share modeling and symmetry
premises; prior exposure to some results prevents a claim of complete
blinding. For rows 4 and 14, the terminal matrix $J$ comes from the
saved derivative reduction, not the final global matrix. Only the
row-10 direct path is used here. The general step-shrinking rounding
counterexample remains valid; the six selected paths use audited normal
accepted calls. Rows 11, 15 and 16 retain Arb-only certification;
the second-route certificates for rows 11 and 16 remain on hold.

Exact artifact locations, input and source hashes, settings, retained
unsuccessful attempts, comparison records, and the provenance of the
mathematical reading are documented in the accompanying
\texttt{VERIFICATION-RECORD.md}.

\section{Model and quarter-period symmetry}
\label{sec:model}
All three masses and the gravitational constant are one. Positions and
velocities are $q_i,p_i\in\R^2$, with
\begin{equation}
 \dot q_i=p_i,\qquad
 \dot p_i=\sum_{k\ne i}\frac{q_k-q_i}{\lVert q_k-q_i\rVert_2^3}.
 \label{eq:newton}
\end{equation}
For $x=(a,b,T)$ we impose
\begin{equation}
 \begin{split}
 q_1(0)&=(-1,0),\quad q_2(0)=(1,0),\quad q_3(0)=(0,0),\\
 p_1(0)&=p_2(0)=(a,b),\qquad p_3(0)=-2(a,b).
 \end{split}
 \label{eq:initial}
\end{equation}
The reduced map is
\begin{equation}
 F(a,b,T)=\bigl(p_{1x},p_{1y},p_{2x}\bigr)(T/4).
 \label{eq:F}
\end{equation}
It is defined on any parameter box for which the required flows exist
without collision through all relevant terminal times.

Each input row has fields $(a_0,b_0,T_0,T^*)$. We integrate through a
quarter of the third field $T$, rather than the fourth, scale-invariant
period $T^*$. The decimal strings are exact rational centers, not assumed
error bounds on the unknown exact solution.

\begin{lemma}[Three terminal components suffice]
\label{lem:symmetry}
Suppose a solution of \eqref{eq:newton}--\eqref{eq:initial} is
collision-free on the closed interval $[0,\tau]$, with $\tau>0$. If
\begin{equation}
 p_{1x}(\tau)=p_{1y}(\tau)=p_{2x}(\tau)=0,
 \qquad q_{2x}(\tau)-q_{3x}(\tau)\ne0,
 \label{eq:terminal}
\end{equation}
then it extends to a collision-free solution on $\R$ with period
$4\tau$. A common lower bound $d>0$ for all three pairwise distances on
$[0,\tau]$ holds on the full extension. If $(a,b)\ne(0,0)$, the solution
is nonconstant.
\end{lemma}
\begin{proof}
Linear momentum $P=\sum_i p_i$ and scalar angular momentum
$L=\sum_i q_i\times p_i$ vanish initially. Pairwise cancellation of
forces gives $\dot P=0$, and the central-force identity gives $\dot L=0$.
The center of mass remains zero, so $q_3=-q_1-q_2$ for the whole
interval. At $\tau$, the three zero components in \eqref{eq:terminal}
and $P=0$ imply $p_{3x}=0$ and $p_{3y}=-p_{2y}$. Thus
$p_{ix}(\tau)=0$ for all $i$ and $p_{1y}(\tau)=0$, so all
$q_{iy}p_{ix}$ terms in the cross products vanish. Hence
\[
 0=L=(q_{2x}-q_{3x})p_{2y}.
\]
The nonzero factor in \eqref{eq:terminal} forces $p_{2y}=p_{3y}=0$.
Thus all velocities vanish at $\tau$.

Time reversal about this brake instant defines
$q(\tau+s)=q(\tau-s)$ for $0\le s\le\tau$. It joins with zero velocity
and solves the same second-order equation. Uniqueness on the
collision-free domain makes it the continuation of the original
solution. Its state at $2\tau$ is $(q(0),-p(0))$.

Define the involution on triples of planar vectors by
\[
 S(z_1,z_2,z_3)=(-z_2,-z_1,-z_3).
\]
It is the composition of a swap of the first two equal masses and a
rotation through $180$ degrees, and is a symmetry of the equations.
For \eqref{eq:initial}, $Sq(0)=q(0)$ and $Sp(0)=-p(0)$. Define the
second half of the solution by
\[
 (q(2\tau+s),p(2\tau+s))=S(q(s),p(s)),\qquad0\le s\le2\tau.
\]
The initial state of this half equals the state already reached at
$2\tau$, so uniqueness again gives a smooth continuation. At $4\tau$
the state is $S(q(0),-p(0))=(q(0),p(0))$. Periodic continuation defines
the solution for every real time. Time reversal and $S$ preserve the
set of pairwise distances, so the common lower bound persists. A
nonzero initial velocity proves nonconstancy.
\end{proof}

The nondegeneracy test is essential: three zero velocity components
without the nonzero factor in \eqref{eq:terminal} would not force all
velocities to vanish. In center-of-mass coordinates the tested quantity
is $q_{2x}-q_{3x}=q_{1x}+2q_{2x}$. Under the hypotheses of
Lemma~\ref{lem:symmetry}, a zero of $F$ determines a collision-free
periodic solution.

\section{Validated integration and the finite-dimensional test}
\label{sec:method}
We describe sufficient conditions used by the computations. Every
interval operation in this section is an inclusion operation. Matrix
norms are induced infinity norms, and boxes are closed rectangular
interval sets. A positive distance bound is checked on a whole time
slab, not on sampled trajectory points.

\subsection{An eight-dimensional center-of-mass system}
Write
\[
 y=(x_1,y_1,x_2,y_2,u_1,w_1,u_2,w_2),
 \qquad q_3=-q_1-q_2,\quad p_3=-p_1-p_2.
\]
With $d_{12}=q_2-q_1$, $d_{13}=-2q_1-q_2$ and
$d_{23}=-q_1-2q_2$, set $g(d)=d/(d\cdot d)^{3/2}$. Then
\[
 \dot p_1=g(d_{12})+g(d_{13}),\qquad
 \dot p_2=-g(d_{12})+g(d_{23}).
\]
Together with $\dot q_i=p_i$, these equations define $y'=f(y)$ and
recover \eqref{eq:newton}. For
\[
 B(d)=(d\cdot d)^{-3/2}I_2-3(d\cdot d)^{-5/2}dd^{\mathsf T},
\]
the Jacobian has block form
\begin{equation}
 A=Df=\begin{pmatrix}0&I_4\\H&0\end{pmatrix},\qquad
 H=\begin{pmatrix}
 -B_{12}-2B_{13}& B_{12}-B_{13}\\
 B_{12}-B_{23}&-B_{12}-2B_{23}
 \end{pmatrix},
 \label{eq:jacobian}
\end{equation}
where $B_{ij}=B(d_{ij})$. Every evaluation requires a strictly positive
lower bound for each $d_{ij}\cdot d_{ij}$. Polynomial squaring and
positive square roots are evaluated with their respective inclusion
rules.

\subsection{Existence and variational bounds on one slab}
Let $Y$ enclose all starting states of a step of rational length $h>0$.
Choose a collision-free box $B$ and a fixed positive diagonal matrix
$D$ for which
\begin{equation}
 Y+[0,h]f(B)\subseteq B,\qquad
 \alpha\ge\sup_{z\in B}\norm{D^{-1}A(z)D},\qquad h\alpha<\tfrac14.
 \label{eq:picard}
\end{equation}
For any starting state in $Y$, the Picard operator maps continuous
$B$-valued paths into themselves. In the weighted sup norm
$\sup_t\norm{D^{-1}y(t)}$ it is a contraction with constant at most
$h\alpha$. Thus a unique solution exists on the entire slab and stays
in $B$.

Let $M$ be the derivative of the step flow, with $M(0)=I_8$. Its
equation is $M'=A(y)M$. The integral equation, in the same weighted
norm, gives
\begin{equation}
 \sup_{0\le t\le h}\norm{D^{-1}M(t)D}
 \le\frac1{1-h\alpha},\qquad
 \norm{D^{-1}(M(t)-I_8)D}\le E:=\frac{h\alpha}{1-h\alpha}.
 \label{eq:variation}
\end{equation}
In particular, a componentwise enclosure for every such $M(t)$ is
\begin{equation}
 \mathcal M_{ij}=\delta_{ij}+[-E D_{ii}/D_{jj},E D_{ii}/D_{jj}].
 \label{eq:Mbox}
\end{equation}

The unweighted choice is $D=I_8$. The weighted implementation uses
$D=\operatorname{diag}(I_4,wI_4)$ with a positive point dyadic $w$ fixed
throughout each slab. By \eqref{eq:jacobian}, a certified bound
$R_H\ge\sup_B\norm{H}$ supplies
\begin{equation}
 \alpha=\max(w,R_H/w).
 \label{eq:weight}
\end{equation}
When $R_H>0$, one may choose a dyadic $w\ge\sqrt{R_H}$; for $R_H=0$,
one may take $w=1$. The method also permits $w<1$. This changes the
bound used for a local proof, not the vector field. Since $D$ is fixed
within the slab, there is no derivative of $D$ in the variational
equation. Taylor remainders and terminal contractions are tested in
the original coordinates.

\subsection{Taylor remainders and propagation}
Normalized Taylor coefficients $c_k$ of the state satisfy
\begin{equation}
 (k+1)c_{k+1}=[s^k]f\left(\sum_{j=0}^k c_js^j\right).
 \label{eq:taylor}
\end{equation}
If $A(y(s))=\sum A_ks^k$, the variational coefficients satisfy
$(k+1)M_{k+1}=\sum_{j=0}^k A_jM_{k-j}$. The recurrences are evaluated
with outward interval operations.

For an order $N$ step from a point $c$, evaluate the state polynomial
through degree $N-1$ at $c$ and enclose its degree $N$ remainder using
the coefficient evaluated on $B$. Thus an enclosure $C$ of $\phi_h(c)$
is
\begin{equation}
 C=\sum_{k=0}^{N-1}c_k(c)h^k+c_N(B)h^N.
 \label{eq:remainder}
\end{equation}
This follows componentwise from Taylor's theorem, since all intermediate
states belong to $B$. Apply the same argument to the joint system
$(y,M)'=(f(y),A(y)M)$. Its initial set is $(Y,I_8)$ and its remainder
set is $(B,\mathcal M)$. This gives an interval matrix $\mathcal D$
containing $D\phi_h(z)$ for every $z\in Y$. The local remainder
tolerance is an additional acceptance criterion; the existence and
inclusion arguments are \eqref{eq:picard}--\eqref{eq:remainder}.

To control wrapping, represent the starting set by $c+Qb$, where $Q$
is an invertible point matrix and $b$ is a zero-centered box; choose $Y$
as a rectangular enclosure. Track the center-parameter solution
separately in $c+Qp$, maintaining $p\subseteq b$. The mean-value
formula on the convex box $Y$ gives
\[
 \phi_h(c+Qb)\subseteq C+\mathcal DQb.
\]
Choose a new point center $c_+\in C$ and point basis $Q_+$. An interval
inverse of $Q_+$ yields zero-centered enclosures
\begin{equation}
 \begin{split}
 b_+&\supseteq Q_+^{-1}(C-c_++\mathcal DQb),\\
 p_+&\supseteq Q_+^{-1}(C-c_++\mathcal DQp).
 \end{split}
 \label{eq:propagation}
\end{equation}
A floating-point QR factorization may propose a basis, but its
orthogonality is not an assumption: invertibility and the inverse
enclosure are checked for the chosen point matrix.

The $8\times2$ sensitivity $W=\partial y/\partial(a,b)$ starts with
four zero position rows and velocity rows $(1,0),(0,1),(1,0),(0,1)$.
If $W\in Q R_W$, update
$R_{W,+}\supseteq Q_+^{-1}\mathcal DQ R_W$. This encloses the
derivative for all parameters in the box, not only its center.

\subsection{The variable terminal time}
The integrator reaches an exact dyadic time $t_0$ near $T_0/4$.
Every step length is rational and positive, and the saved lengths sum
exactly to $t_0$. For $X=x_0+[-r,r]^3$, choose
\begin{equation}
 \Delta\ge\abs{T_0/4-t_0}+r/4.
 \label{eq:Delta}
\end{equation}
An endpoint box $B_e$ is accepted when
\begin{equation}
 Y_e+[-\Delta,\Delta]f(B_e)\subseteq B_e,\qquad
 \Delta\sup_{B_e}\norm{D^{-1}AD}<\tfrac14,
 \label{eq:endpicard}
\end{equation}
where $Y_e$ encloses the parameter-box states at $t_0$. The positive-
and negative-time Picard arguments then cover all required terminal
times. If $\delta_0$ encloses the exact rational $T_0/4-t_0$, the
center solution at $T_0/4$ belongs to
$c+Qp+\delta_0 f(B_e)$. Its three selected velocities give an enclosure
$F_0$ of $F(x_0)$. Errors from converting the decimal centers to binary
are included in $p$ and $\delta_0$.

For completeness, let $W_0$ enclose $W(t_0)$, choose
$\alpha_e\ge\sup_{B_e}\norm{D^{-1}AD}$, and put
$E_e=\Delta\alpha_e/(1-\Delta\alpha_e)$. If
$n\ge\norm{D^{-1}W_0}$, an enclosure throughout the small terminal
interval is
\[
 W_e=W_0+\{V:\abs{V_{ij}}\le D_{ii}E_en\}.
\]
An endpoint derivative enclosure is consequently
$W_0+[-\Delta,\Delta]A(B_e)W_e$. Selecting its velocity rows
$(1x,1y,2x)$ gives the $a,b$ columns of an interval matrix $J$.
The $T$ column consists of the same components of $f(B_e)/4$.
The chain rule therefore gives $DF(X)\subseteq J$.

\subsection{Krawczyk contraction and collision margins}
Choose an invertible point matrix $C_0\in\R^{3\times3}$, put
$Z=I_3-C_0J$, and form the centered Krawczyk enclosure
\begin{equation}
 K_c=-C_0F_0+Z[-r,r]^3.
 \label{eq:Krawczyk}
\end{equation}
The finite-dimensional tests are
\begin{equation}
 K_c\subset\operatorname{int}[-r,r]^3,\qquad \norm{Z}<1.
 \label{eq:Kconditions}
\end{equation}

\begin{proposition}
\label{prop:contraction}
Suppose $F$ is continuously differentiable on a neighborhood of $X$,
$F(x_0)\in F_0$, $DF(X)\subseteq J$, $C_0$ is invertible, and
\eqref{eq:Kconditions} holds. Then $F$ has exactly one zero in $X$.
\end{proposition}
\begin{proof}
For $G(x)=x-C_0F(x)$, the integral mean-value formula shows
$G(X)\subseteq x_0+K_c\subseteq\operatorname{int}X$.
Also $DG(X)\subseteq Z$, so $G$ is a contraction on the complete
space $X$. Banach's theorem gives exactly one fixed point in $X$.
Invertibility of $C_0$ identifies fixed points of $G$ with zeros of $F$.
\end{proof}

The validated flow on collision-free slabs provides the smoothness
needed here.

Every slab and the terminal box must have a common positive rational
lower bound $d_2$ for all three squared pairwise distances. A displayed
distance bound $d$ is rounded down so that $d^2\le d_2$. We also require
that the interval for $q_{1x}+2q_{2x}$ on $B_e$ exclude zero, that the
parameter box have $T>0$, and that its initial velocity exclude
$(0,0)$. At the zero supplied by Proposition~\ref{prop:contraction},
Lemma~\ref{lem:symmetry} then gives the claimed nonconstant periodic
solution and the full-time collision margin.

\section{Separate implementations and identification of the zeros}
\label{sec:implementations}
Implementation A uses Python 3.14.7, python-flint 0.9.0 and FLINT 3.6.0
with Arb real balls. Implementation B is a separate C++ program using
MPFR 4.2.2 and GMP 6.3.0 with lower and upper interval endpoints. Both
implement the Taylor method of Section~\ref{sec:method}, with separate
interval and basis handling. Implementation C uses CAPD 6.1.0 with
multiprecision enabled, MPFR 4.2.0 and GMP 6.2.1. Its C1-Lohner
integration, normalized clock and derivative reduction are described
below. A versus B and A versus C provide two flow routes per doubly
certified row, using three codebases overall.

\begin{table}[htbp]
\centering\small
\begin{tabular}{@{}lcccc@{}}
\toprule
Run family & Bits & Order & Parameter radius & Tolerance\\
\midrule
A, rows 2, 3 & 320 & 24 & $2^{-120}$ & $2^{-170}$\\
A, other successes & 448 & 48 & $2^{-160}$ & $2^{-250}$\\
B, all successes & 384 & 28 & $2^{-100}$ & $2^{-180}$\\
C, rows 4, 13, 14: $C^0$ & 320 & 24 & Point enclosure & $2^{-220}$\\
C, rows 4, 13, 14: $C^1$ & 256 & 24 & $2^{-90}$ or $2^{-100}$ & $2^{-180}$\\
C, rows 5, 6, 10: $C^0$ & 448 & 56 & Point enclosure & $2^{-300}$\\
C, rows 5, 6, 10: $C^1$ & 320 & 48 & $2^{-130}$ or $2^{-140}$ & $2^{-260}$\\
\bottomrule
\end{tabular}
\caption{Selected successful settings. A uses an unweighted Picard
norm on rows 2 and 3 and a fixed diagonal weight on the other rows;
B uses the unweighted norm. C uses CAPD doubleton enclosures. These
run families are not a common-setting benchmark.}
\label{tab:settings}
\end{table}

\subsection{CAPD: a fixed clock for a variable period}
Adjoin the period as a constant ninth state, $\theta=T$, and solve
\begin{equation}
 \frac{d}{ds}(y,\theta)=\bigl((\theta/512)f(y),0\bigr),
 \qquad 0\le s\le128.
 \label{eq:clock}
\end{equation}
For each initial parameter triple the physical time is
$t=\theta s/512$, so the fixed endpoint $s=128$ is exactly $t=T/4$.
The entries $p_i$ remain physical velocities. The affine initial map
$E(a,b,T)$ encodes \eqref{eq:initial} and $\theta=T$; a three-column
doubleton carries the correlated parameter box $E(Y_j)$. Outward
conversion of the exact rational center is included. The center
computation encloses $E(x_{0,j})$ and supplies $F(x_{0,j})$; the box
computation supplies the derivative over every parameter in $Y_j$.

The nonvalidated reference trajectory serves progress and cost
estimation. CAPD's step controller selects the accepted steps; reference
samples neither prescribe a certified partition nor replace the evolving
initial family. In particular, a reference interval count must not be
identified with the number of accepted slabs.

Let $H_k$ be the saved positive interval of step lengths. The audits
establish that every nonfinal $H_k$ is a singleton. Choose its rational
value $h_k$, and choose the final length as
\[
 h_{m-1}=128-\sum_{k<m-1}h_k.
\]
The final choice is positive and belongs to the final $H_{m-1}$; it need
not equal its left endpoint. Put $s_0=0$ and $s_{k+1}=s_k+h_k$.
These are exact rational times common to all parameters, and $s_m=128$.
Membership of a length in $H_k$ is used only after a uniform local
inclusion over that interval has been established.

For the selected normal calls, let $\Pi_k$ enclose the Taylor polynomial
on $[0,\sup H_k]$, $M_k$ be its remainder margin, $\mathcal E_k$ the
candidate enclosure, and $R_k$ the remainder evaluated on that enclosure.
The state and local variational components satisfy
\[
 R_k\subset\operatorname{int} M_k,\qquad
 \Pi_k+M_k\subseteq\mathcal E_k,\qquad
 0\in\operatorname{int}M_k.
\]
The time-zero input state box and variational identity matrix are
contained in the corresponding polynomial ranges and enclosures. Strictly
positive squared distances on the state tube make the extended ODE
smooth there. Its Taylor formula and a first-exit argument give enclosure
on $[0,\sup H_k]$, uniformly for every $h\in H_k$. The local $9\times9$
matrix $L_k$ encloses the flow derivative for the input box and every
such $h$, including the variational remainder. The fixed-source audit
connects these polynomial and remainder records to CAPD's Taylor/AD
calculation; the separate integer checks verify their inclusions.

To connect these local statements to the original parameters, write
$u=(a,b,T)\in Y_j$, $S=DE$, and, for the normalized flow $\Phi$,
\[
 Y_k(u)=\Phi_{s_k}(E(u)),\qquad V_k(u)=D_uY_k(u).
\]
At each selected time, both the cached axis box $B_k$ and the geometric
doubleton $\mathcal R_k$ contain $Y_k(u)$. They need not contain each
other. The set update preserves this reachable-family assertion; no
assertion that the saved output contains the image of every fictitious
point of $B_k$ is required. The matrix family carried by the set update
contains a true matrix $A_k(u)$ with
\[
 A_k(u)S=V_k(u),\qquad A_0=I_9.
\]
The actual derivative obeys the chain rule at the same $u,s_k,h_k$.
Thus the local product and the saved global $9\times3$ enclosure
$G_{k+1}$ both contain $V_{k+1}(u)$. This simultaneous state and tangent
induction, starting with the original affine family, justifies their
intersection. It does not reset a state or a derivative at a reference
sample. The center calculation has the same time-coverage construction
but starts with the center enclosure only, independently of the command
line's radius label.

\subsection{Eliminating the connection variables}
Let $\phi_k(z_k,T)-z_{k+1}=0$ be the eight-state connection equation
between consecutive nodes. With the initial node prescribed by the
initial map, the Jacobian with respect to the subsequent nodes
$z_1,\ldots,z_m$ is block lower triangular, with diagonal blocks $-I$.
Forward elimination of these variables transports the derivative with
respect to $(a,b,T)$, retaining the parameter dependence of the prescribed
initial node. Using the ninth constant state includes the direct period
dependence:
\begin{equation}
 D_0\supseteq DE,\qquad
 D_{k+1}=\operatorname{out}(L_kD_k)\cap G_{k+1}.
 \label{eq:elimination}
\end{equation}
Here $D_k$ is an interval enclosure of $V_k(u)$ for every original
parameter $u$. Interval products are rounded outward onto the
$2^{-240}$ grid for rows 4, 13 and 14, and the $2^{-512}$ grid for
rows 5, 6 and 10. Intersecting each propagated enclosure with CAPD's
global parameter-derivative enclosure preserves inclusion, because both
contain the true derivative for the same initial family and time node.
The final rows $(p_{1x},p_{1y},p_{2x})$ give a $3\times3$ enclosure $J$
of $DF(Y_j)$.

The proof uses coverage of every true trajectory starting in $E(Y_j)$.
Consequently the reduced derivative is a derivative enclosure on the
whole original parameter box. Uniqueness of an isolated root of an
expanded shooting system alone would not give that coverage. The
reduction uses local and global enclosures from the same CAPD flow; they
are not two independent integrations.

The terminal MPFR calculation uses 512 bits for rows 4, 13 and 14,
and 1024 bits for rows 5, 6 and 10, and applies
\eqref{eq:Krawczyk}--\eqref{eq:Kconditions}. The audits separately
propagate a reduced enclosure using a $2^{-512}$ grid and verify that
the saved terminal derivative contains it. For rows 4, 13 and 14 this is
an inclusion, not numerical equality of the old and audited reductions.
The terminal $J$ is linked to the velocity rows of the saved reduction;
the terminal $F_0$ and nondegeneracy factor are linked to the center and
box endpoints, respectively. Exact rational checks of these saved
terminal inputs establish invertibility, strict inclusion, contraction
and the nonzero factor. The six local-inclusion audits supply the
additional connection to the flow described above; terminal algebra
alone would not supply it.

\subsection{Identifying the two zeros for each row}
\label{sec:samezero}
\begin{proposition}[Identification by nested boxes]
\label{prop:same}
Suppose the same reduced equation $F$ has one certified zero in each
of two boxes $X\subseteq Y$, and has exactly one zero in $Y$. Then the
two zeros coincide.
\end{proposition}
\begin{proof}
The zero in $X$ belongs to $Y$ and must equal its unique zero.
\end{proof}

For all ten rows, the saved comparisons establish equal exact rational
centers. The radii in Theorem~\ref{thm:result} give $X_j\subseteq Y_j$.
Proposition~\ref{prop:same} therefore identifies the two zeros. A and B
cover the variable endpoint by \eqref{eq:Delta}; C covers it by
\eqref{eq:clock}. All represent the same map $F$. Overlap of endpoint
or residual intervals is a useful consistency check, but nesting and
uniqueness are the reason for the same-solution conclusion.

Rows 1, 2, 3 and 12 have B's unweighted Taylor route. Rows 4, 5, 6,
10, 13 and 14 have C's CAPD route. These ten second-route certificates
do not rely on A's weighted integration bounds. Rows 11, 15 and 16
remain certified by A alone in the adopted evidence; partial CAPD
records on rows 11 and 16 do not add a second certificate.

\section{Results by table row}
\label{sec:results}
The present result set concerns rows 1--20 of the 4,860-row input.
Table~\ref{tab:rows} distinguishes certification, unsuccessful
terminal tests, partial integration and unstarted runs. Counts describe
coverage over retained attempts with varying settings, rather than a
success rate under a single fixed protocol.

\begingroup
\small
\setlength{\LTcapwidth}{\textwidth}
\setlength{\LTpre}{4pt}
\begin{longtable}{@{}rcccc@{}}
\toprule
Row & A (Arb) & B (MPFR) & C (CAPD) & Two routes, same zero\\
\midrule
\endfirsthead
\toprule
Row & A (Arb) & B (MPFR) & C (CAPD) & Two routes, same zero\\
\midrule
\endhead
1 & $\Cert$ & $\Cert$ & $\Notrun$ & Yes\\
2 & $\Cert$ & $\Cert$ & $\Notrun$ & Yes\\
3 & $\Cert$ & $\Cert$ & $\Notrun$ & Yes\\
4 & $\Cert$ & $\Partial$ & $\Cert$ & Yes\\
5 & $\Cert$ & $\Partial$ & $\Cert$ & Yes\\
6 & $\Cert$ & $\Partial$ & $\Cert$ & Yes\\
7 & $\Uncert$ & $\Partial$ & $\Partial$ & ---\\
8 & $\Uncert$ & $\Partial$ & $\Notrun^{*}$ & ---\\
9 & $\Uncert$ & $\Partial$ & $\Notrun^{*}$ & ---\\
10 & $\Cert$ & $\Partial$ & $\Cert$ & Yes\\
11 & $\Cert$ & $\Partial$ & $\Hold$ & ---\\
12 & $\Cert$ & $\Cert$ & $\Notrun$ & Yes\\
13 & $\Cert$ & $\Partial$ & $\Cert$ & Yes\\
14 & $\Cert$ & $\Notrun$ & $\Cert$ & Yes\\
15 & $\Cert$ & $\Notrun$ & $\Uncert/\Partial$ & ---\\
16 & $\Cert$ & $\Notrun$ & $\Hold$ & ---\\
17 & $\Uncert$ & $\Notrun$ & $\Partial$ & ---\\
18 & $\Uncert$ & $\Notrun$ & $\Notrun^{*}$ & ---\\
19 & $\Uncert$ & $\Notrun$ & $\Notrun^{*}$ & ---\\
20 & $\Uncert$ & $\Notrun$ & $\Partial$ & ---\\
\midrule
Certified rows & 13 & 4 & 6 & 10\\
\bottomrule
\caption{Status over retained attempts. $\Cert$: certified;
$\Uncert$: attempted but uncertified; $\Partial$: incomplete flow;
$\Notrun$: no full-flow attempt. $\Notrun^{*}$: local pilot tests only.
$\Hold$: second certificate on hold, with incomplete audit evidence.
For C on row 15, $\Uncert/\Partial$ retains terminal failures and
incomplete integrations; it does not imply a second certificate. A dash makes no
same-solution assertion.}
\label{tab:rows}
\end{longtable}

For the distance columns, Table~\ref{tab:distances} gives downward-rounded
linear bounds $d_A$, Table~\ref{tab:second} displays approximations of
$\sqrt{d_{2,B}}$, and Table~\ref{tab:capd} displays approximations of
the squared bounds $d_{2,C}$.

\begin{longtable}{@{}rrrl@{}}
\toprule
Row & A steps & A radius & Full-time distance lower bound $d_A$\\
\midrule
1 & 1,885 & $2^{-160}$ & 0.287841\\
2 & 39,860 & $2^{-120}$ & 0.014202\\
3 & 47,639 & $2^{-120}$ & 0.015223\\
4 & 8,047 & $2^{-160}$ & 0.000363\\
5 & 10,809 & $2^{-160}$ & 0.002044\\
6 & 10,922 & $2^{-160}$ & 0.002188\\
10 & 12,963 & $2^{-160}$ & 0.000286\\
11 & 13,080 & $2^{-160}$ & 0.000266\\
12 & 5,670 & $2^{-160}$ & 0.287755\\
13 & 8,813 & $2^{-160}$ & 0.002304\\
14 & 8,902 & $2^{-160}$ & 0.001999\\
15 & 14,991 & $2^{-160}$ & 0.000007\\
16 & 13,321 & $2^{-160}$ & 0.000012\\
\bottomrule
\caption{Selected A certificates. Every decimal in the last column
is a downward-rounded bound with $d_A^2\le d_2$. The slab bound holds
through the quarter-period and transfers to all time by
Lemma~\ref{lem:symmetry}. Different successful certificates for the
same row can give different valid lower bounds.}
\label{tab:distances}
\end{longtable}

\begin{longtable}{@{}rrccc@{}}
\toprule
Row & B steps & B radius & Approximate $\sqrt{d_{2,B}}$ & Approximate contraction bound\\
\midrule
1 & 2,326 & $2^{-100}$ & 0.288723577 & $1.48\times10^{-18}$\\
2 & 31,139 & $2^{-100}$ & 0.014202810 & $3.58\times10^{-3}$\\
3 & 38,951 & $2^{-100}$ & 0.015223399 & $6.31\times10^{-3}$\\
12 & 6,995 & $2^{-100}$ & 0.288635565 & $1.70\times10^{-6}$\\
\bottomrule
\caption{B certificates on four of the ten doubly certified solutions. The
last two columns are display approximations of saved quantities;
they are not rounded bounds used in the proof. Strict positivity,
strict inclusion and contraction are decided on the saved binary
intervals and their exact rational interpretation.}
\label{tab:second}
\end{longtable}

\begin{longtable}{@{}rrrccc@{}}
\toprule
Row & C $C^0$ slabs & C $C^1$ slabs & Radius & Approximate $d_{2,C}$ & $q_C$ (approx.)\\
\midrule
4 & 37,024 & 12,325 & $2^{-90}$ & $1.323016679\times10^{-7}$ & 0.1430438542\\
5 & 4,880 & 4,868 & $2^{-140}$ & $4.177760495\times10^{-6}$ & 0.0002962666\\
6 & 4,884 & 4,871 & $2^{-130}$ & $4.801121064\times10^{-6}$ & 0.1317961327\\
10 & 5,752 & 5,849 & $2^{-140}$ & $8.134207033\times10^{-8}$ & 0.0003566111\\
13 & 36,489 & 12,134 & $2^{-90}$ & $5.315061159\times10^{-6}$ & 0.0298201405\\
14 & 36,875 & 12,264 & $2^{-100}$ & $4.028534374\times10^{-6}$ & 0.0102606219\\
\bottomrule
\caption{C certificates. The squared-distance bounds and contraction
bounds $q_C=\norm{I-C_0J}$ are display approximations; the contraction
column transcribes the saved terminal field \texttt{q\_upper}. This is
distinct from the full Krawczyk inclusion ratio defined in
Section~\ref{sec:certificate}. The saved exact intervals decide the
inequalities. Strict Krawczyk inclusion and the
nonzero endpoint factor also pass in all six cases. $C^0$ counts
include reused center evaluations, which are independent of box radius.}
\label{tab:capd}
\end{longtable}
\endgroup

A attempts every row from 1 to 20. Its seven uncertified rows are
7, 8, 9 and 17--20. Minimum-step stops and time limits give no
nonexistence conclusion. A positive distance bound on a completed
prefix does not bound an unintegrated portion of the quarter-period.

For B, rows 4--11 and 13 remain partial, and rows 14--20 have no
full-flow attempt. Historical C tests on rows 5, 6, 10, 11 and 16 failed
at radii $2^{-90}$ and $2^{-100}$. The smaller-box successes for rows
5, 6 and 10 in Table~\ref{tab:capd} supersede those failures only for
the stated boxes and selected paths.

For rows 11 and 16, the available second-route evidence does not establish
all required flow and terminal connections. Neither partial result is
adopted as a second certificate.

Later full box attempts on rows 7, 15, 17 and 20 stopped at normalized
times approximately 50.1671, 72.2639, 41.0290 and 41.3307 of 128;
their full center runs did not start. Rows 8, 9, 18 and 19 have local
pilots only. The earlier row-15 $2^{-90}$ terminal failure and
$2^{-100}$ stop after 26,140 slabs at time 108.646304 are retained.
None of these incomplete or unsuccessful tests proves nonexistence.
In particular, rows 11, 15 and 16 still have A's existence certificates.

\subsection{Measured work and recorded cost}
The CAPD phase times reported below are observations from the earlier
order-24 runs for rows 4, 13 and 14; they do not include the later
certifications of rows 5, 6 and 10.
The selected B runs take approximately 241, 1,963, 2,338 and 507 seconds
for rows 1, 2, 3 and 12. For rows 1 and 12 these values are taken from
the certificates' elapsed-time fields; for rows 2 and 3 they are rounded
worker wall times. C's successful $C^1$ native runs take about
3,704, 3,093 and 2,831 seconds for rows 4, 13 and 14, respectively.
Each requires a center evaluation as well: about 3,758, 3,735 and
3,732 seconds, with reuse across radii at the same exact center.
C's phase times exclude compilation, record checks and transfer.
Different boxes, precisions and machines preclude a controlled speed
comparison between implementations.

The historical eight-row C batch covering rows 4, 5, 6, 10, 11, 14, 15 and 16 has
22 completed integrations, one retained partial run and three platform
preemptions. It uses one physical core and 1\,GiB per flow process, with
at most eight concurrent cores. The largest recorded peak resident
memory among the completed native integrations is 63,672\,KiB. The
batch's conservative recorded cost estimate is US\$2.393, including
unsuccessful attempts and a US\$0.50 transfer/storage allowance. It is
an estimate, not an observed invoice. The row-13 full-box attempt and
smaller-box certification have separate recorded estimates of
US\$0.340 and US\$0.225; the latter reuses the earlier center evaluation.
These amounts exclude earlier development, local pilots, the later
row-5/6/10 smaller-box runs and subsequent method audits. They measure
the indicated historical batches, not the total cost of this article.
Shared credit-ledger balances and audit expenditure are different
accounting units and are not added to these figures. Local-pilot
extrapolations are not measurements of full-orbit speed or cost.

The table contains 4,840 further rows, 21--4,860, outside A's attempted
range and outside the present claim. Extending a run or improving an
enclosure may change an uncertified status. No conclusion about the
full database follows from this finite prefix.

\begin{proof}[Proof of Theorem~\ref{thm:result}]
For A and B, the full sequence of accepted slabs satisfies the Picard
self-map and contraction tests, positive distance tests, and state and
variational Taylor remainder bounds. Verified basis changes preserve
the enclosure inductively. The final time correction gives $F_0$ and
$J$. For C, the audited uniform local inclusions and reachable-family/tangent
induction cover the original initial family at common rational times;
connection elimination encloses its parameter derivative. Thus each route supplies the inclusions required
in Proposition~\ref{prop:contraction}.
The saved terminal tests pass exact rational checks of invertibility,
strict interior inclusion and contraction. The positive distance bound,
nonzero terminal factor, positive period and nonzero initial velocity
complete the hypotheses of Lemma~\ref{lem:symmetry}. Thus every stated
box contains its unique reduced zero and supplies the periodic solution
and collision margin. For the ten doubly certified rows, the saved
exact-center and box-inclusion comparisons, together with
Proposition~\ref{prop:same}, identify the zeros.
\end{proof}

\section{Certificate format and levels of checking}
\label{sec:certificate}
A and B certificates are JSON objects containing the input hash, row
number, exact decimal center, radius exponent, precision, Taylor order,
norm choice, remainder tolerance, complete rational step list, and its
exact total time. Each certificate also contains the final point center and basis,
the point and parameter-box enclosures, the sensitivity enclosure,
$F_0,J,C_0,K_c,Z,B_e$, the common squared distance bound, and the
nondegeneracy interval. An interval ball is saved as
\begin{equation}
 m2^a\ \mathbin{\pm}\ n2^b,\qquad m,a,n,b\in\mathbb Z,\quad n\ge0,
 \label{eq:serialization}
\end{equation}
with integers encoded as strings where necessary. Matrix entries are
stored row by row. This preserves exact binary endpoints independently
of a decimal display. Point matrices have radius zero.

For C, the contraction bound $q_C$ and the full inclusion ratio
$\rho_C=\max_i\sup_{u\in(K_c)_i}|u|/r$ are separate quantities.
The saved fields are \nolinkurl{terminal.mpfr.q_upper} and
\nolinkurl{terminal.fraction.max_inclusion_ratio_float}, respectively.
The latter is a display approximation; exact interval and rational tests
decide strict inclusion. For example, row 13 has approximately
$q_C=0.02982014046783021$ and $\rho_C=0.02982014046810953$.
Their agreement at the shorter precision of Table~\ref{tab:capd} does
not identify the two quantities: $\rho_C$ includes the center residual
term in \eqref{eq:Krawczyk}.

The terminal certificate serializes a finite algebraic implication
conditional on its flow and derivative enclosures. A success flag alone
does not establish those analytic premises. Three additional kinds of
evidence must be distinguished: source examination, saved local-inclusion
and reduction checks, and a fresh validated integration replay. For C,
the six selected paths have the first two kinds, including all final
accepted-call records and the uniform-$H$ argument in Section~4.1.
They do not have independent regeneration of Taylor coefficients or
new ODE reintegration from these audits. The trust boundary is stated
once in Section~\ref{sec:assurance}.

Separate full-flow computations by B and C cover the ten specified
rows. C's separately written saved-record checks and source audits
connect its local inclusions to its terminal enclosures; a terminal
rational checker by itself cannot do so.

An unsuccessful Picard, tolerance or terminal inclusion test is
inconclusive. In particular, failure of strict Krawczyk inclusion
does not prove that no zero exists. The documented rounding
counterexample in CAPD's general step-shrinking formula is retained;
the result here relies on the selected normal accepted calls, not on a
general claim that every step-selection branch is correct.

\section{Relation to earlier work and remaining questions}
\label{sec:prior}
Hristov, Hristova, Puzynina, Sharipov and Tukhliev
\cite{Hristov} supply the numerical catalogue and its quarter-period
search formulation; their Section 4.2 imposes the vanishing of all six
velocity components. Their published Euler-coordinate table
\cite{Data} supplies every center in this work. Montgomery's treatment
of brake orbits \cite[Section 3]{Montgomery} explains the role of time
reversal, equal-mass permutations and Euler configurations.
Lemma~\ref{lem:symmetry} states and proves the sufficient direction
used here, with the additional nondegeneracy check that makes three
velocity components sufficient.

Computer-assisted existence proofs in the three-body problem have
substantial precedents. Kapela and Zgliczy\'nski
\cite[Sections 2--3]{KapelaZgliczynski} combine rigorous flow and
derivative integration with interval Newton and Krawczyk methods in
their proof of the figure eight and its local uniqueness. Kapela and
Sim\'o \cite{KapelaSimo} develop computer-assisted existence and KAM
stability methods and apply them to the figure eight and selected
rotating eights. We use no stability conclusion from that work for
the brake orbits studied here. Johansson's Arb library
\cite{Johansson} supplies rigorous ball arithmetic through the
python-flint interface \cite{PythonFlint}; MPFR supplies directed,
correctly rounded scalar operations for B and C
\cite{MPFR}. These contributions provide the objects, symmetry ideas,
existence-proof methods and arithmetic on which this work depends.

Earlier existence results for periodic brake orbits include the
isosceles solutions of Moeckel, Montgomery and Venturelli
\cite[Theorem 2 and Section 5]{BrakeSyzygy}. Those solutions encounter
binary collisions and are distinct from the collision-free objects
certified here. The numerical collision-free searches of Li and Liao
\cite{LiLiao} and Hristov et al.\ \cite{Revisited,Hristov} supply the
catalogue context. CAPD's rigorous flow and variational algorithms
\cite{CAPD} supply C's numerical foundation.

\subsection{Correspondence with earlier numerical catalogues}
Among the thirteen Arb-certified entries in Theorem~\ref{thm:result},
ten---1, 2, 3, 4, 5, 6, 10, 12, 13 and 14---have certificates from two
implementations for the same reduced zero. These entries belong to the
central-symmetry free-fall catalogue of Hristov et al.\
\cite{Hristov,Data}. Rows 1, 2, 3, 5, 6, 13 and 14 numerically
match self-dual entries 1, 2, 3, 8, 9, 27 and 28 of the earlier free-fall
catalogue \cite{Revisited,FreeData}, respectively. Entry 8 is also listed
in Tables 2 and 3 of the earlier paper; entry 9 is identified in its
accompanying online catalogue. These correspondences identify the
published numerical candidates; they are not additional
validated orbit-identification theorems. The thirteen Arb certifications
and the ten same-solution comparisons retain the scope stated in
Theorem~\ref{thm:result}.

\begin{table}[htbp]
\centering\small
\begin{tabular}{@{}cp{.28\linewidth}p{.48\linewidth}@{}}
\toprule
2025 row & Earlier full-table row & Earlier self-dual label\\
\midrule
1 & 3 & self-dual \#1\\
2 & 32 & self-dual \#2\\
3 & 35 & self-dual \#3\\
4 & Not matched & Earlier label unconfirmed\\
5 & 126 & self-dual \#8\\
6 & 129 & self-dual \#9\\
10 & Not matched & Earlier label unconfirmed\\
12 & Not matched & Earlier label unconfirmed\\
13 & 807 & self-dual \#27\\
14 & 810 & self-dual \#28\\
\bottomrule
\end{tabular}
\caption{Numerical correspondence status for the ten rows with two flow
implementations. The earlier table \cite{FreeData} contains 24,582
entries in \texttt{sol\_80.txt} and 236 self-dual entries in
\texttt{sd\_80.txt}; all row numbers are one-based. ``Not matched'' means
no match in these reviewed earlier tables, not a new numerical discovery:
all ten entries are published in the 2025 catalogue \cite{Data}.}
\label{tab:crosswalk}
\end{table}

The comparison uses the paired Euler-coordinate and free-fall tables
published with \cite{Hristov,Data}. This pairing provides the change of
representation; the present comparison does not validate it by a new
flow integration. For each of the seven numerical matches in
Table~\ref{tab:crosswalk}, the four free-fall fields (two initial position
coordinates, period and scale-invariant period) agree to an absolute
difference below $5\times10^{-79}$. Agreement of finite decimals identifies
a numerical catalogue entry, but supplies no error bound on an older
approximation and no proof that it equals a certified root. The rigorous
same-root argument in Section~\ref{sec:samezero} instead concerns two
certified boxes for the same 2025 row.

For the original seven compared rows (1, 2, 3, 4, 12, 13 and 14), we
did not identify earlier rigorous existence proofs in the sources
reviewed through 23 October 2026. A supplement for rows 5, 6 and 10,
reviewed through 25 October 2026, likewise did not identify earlier
rigorous existence proofs for those specified entries. These reviews
distinguish figure-eight and isosceles collision-brake results from the
present candidates. Citation-index access was incomplete, and an
identified citing review and a related full text were not obtained.
These are limits of the reviewed evidence; no claim of priority is made.

\paragraph{Conclusion.}
The retained evidence gives thirteen Arb-certified parameter boxes and
ten same-zero comparisons using a second flow route. The six adopted
CAPD paths are supported by the reachable-family and tangent argument
and their selected-trace audits. Rows 11, 15 and 16 retain Arb-only
certification; the second-route certificates for rows 11 and 16 remain
on hold. These are bounded certification
results for published numerical candidates; no count of geometrically
distinct orbits, least-period result, stability result, catalogue
completeness or priority claim follows.

\clearpage
\paragraph{Acknowledgments.}
The author is grateful to Hristov and coauthors for publishing their
numerical discoveries and accompanying catalogues, and to the developers
of Arb, python-flint, MPFR and CAPD for the software used here. These
published data, methods and software make the present certification work
possible.

\paragraph{Use of AI.}
AI systems assisted with the research planning, prose mathematics,
implementation, record checking, mathematical reading, comparison
and drafting. The author is responsible for the claims and presentation;
AI assistance is not human expert review or formal verification.

\begin{thebibliography}{99}
\bibitem{Hristov}
I. Hristov, R. Hristova, T. Puzynina, Z. Sharipov and Z. Tukhliev,
\emph{Numerical search for three-body periodic free-fall orbits with
central symmetry}, arXiv:2503.00432v1 (2025).
\url{https://arxiv.org/abs/2503.00432v1}.

\bibitem{Data}
I. Hristov et al., accompanying numerical data,
\texttt{euler\_100\_4860.txt}, database of 4,860 orbits.
\url{https://db2.fmi.uni-sofia.bg/3bodysym4860/}.

\bibitem{FreeData}
I. Hristov et al., accompanying data for \emph{Three-body periodic
collisionless equal-mass free-fall orbits revisited}, full and self-dual
catalogues, \texttt{sol\_80.txt} and \texttt{sd\_80.txt}.
\url{https://db2.fmi.uni-sofia.bg/3bodyfree/}.

\bibitem{Montgomery}
R. Montgomery, \emph{Dropping Bodies}, arXiv:2207.02582v1 (2022).
\url{https://arxiv.org/abs/2207.02582v1}.

\bibitem{KapelaZgliczynski}
T. Kapela and P. Zgliczy\'nski,
\emph{An existence of simple choreographies for N-body problem---a
computer assisted proof}, arXiv:math/0304404v1 (2003).
\url{https://arxiv.org/abs/math/0304404v1}.

\bibitem{KapelaSimo}
T. Kapela and C. Sim\'o,
\emph{Rigorous KAM results around arbitrary periodic orbits for
Hamiltonian Systems}, arXiv:1105.3235v1 (2011).
\url{https://arxiv.org/abs/1105.3235v1}.

\bibitem{Johansson}
F. Johansson, \emph{Arb: Efficient Arbitrary-Precision Midpoint-Radius
Interval Arithmetic}, IEEE Transactions on Computers 66 (2017),
1281--1292; arXiv:1611.02831.
\url{https://arxiv.org/abs/1611.02831}.

\bibitem{PythonFlint}
python-flint contributors, version 0.9.0 documentation,
\emph{arb---real numbers} and \emph{arb\_mat---matrices over real numbers}.
\url{https://python-flint.readthedocs.io/en/latest/arb.html};
\url{https://python-flint.readthedocs.io/en/latest/arb_mat.html}.

\bibitem{MPFR}
MPFR contributors, \emph{GNU MPFR 4.2.2}, Sections 1 and 4.4.
\url{https://www.mpfr.org/mpfr-4.2.2/mpfr.html}.

\bibitem{CAPD}
T. Kapela, M. Mrozek, D. Wilczak and P. Zgliczy\'nski,
\emph{CAPD::DynSys: a flexible C++ toolbox for rigorous numerical analysis
of dynamical systems}, Communications in Nonlinear Science and Numerical
Simulation 101 (2021), 105578; arXiv:2010.07097.
\url{https://arxiv.org/abs/2010.07097}.

\bibitem{BrakeSyzygy}
R. Moeckel, R. Montgomery and A. Venturelli,
\emph{From Brake to Syzygy}, Archive for Rational Mechanics and Analysis
204 (2012), 1009--1060; arXiv:1109.1089.
\url{https://arxiv.org/abs/1109.1089}.

\bibitem{LiLiao}
X. Li and S. Liao, \emph{Collisionless periodic orbits in the free-fall
three-body problem}, New Astronomy 70 (2019), 22--26;
arXiv:1805.07980. \url{https://arxiv.org/abs/1805.07980}.

\bibitem{Revisited}
I. Hristov, R. Hristova, V. Dmitra\v{s}inovi\'c and K. Tanikawa,
\emph{Three-body periodic collisionless equal-mass free-fall orbits
revisited}, Celestial Mechanics and Dynamical Astronomy 136 (2024), 7;
arXiv:2308.16159. \url{https://arxiv.org/abs/2308.16159}.
\end{thebibliography}
\end{document}
